\documentclass[11pt,reqno]{amsart}

\usepackage{amssymb}
\usepackage[letterpaper,margin=1.18in]{geometry}
\usepackage[hidelinks]{hyperref}
\hypersetup{
  pdftitle={Rudin's conjecture on squares in arithmetic progressions from abc},
  pdfauthor={Alexey Kravatskiy}
}

\newtheorem{theorem}{Theorem}
\newcommand{\rad}{\operatorname{rad}}
\newcommand{\ii}{\mathrm{i}}
\newcommand{\Z}{\mathbb{Z}}

\title[Rudin's conjecture from \texorpdfstring{$abc$}{abc}]
 {Rudin's conjecture on squares in arithmetic progressions from \texorpdfstring{$abc$}{abc}}
\author{Alexey Kravatskiy}
\address{MIRIAI (Moscow Independent Research Institute of Artificial Intelligence),
Moscow, Russia}
\email{kravatskii.a@miriai.org}
\date{November 27, 2026}
\subjclass[2020]{11B25, 11D45, 11N64}
\keywords{Rudin's conjecture, arithmetic progressions, squares, $abc$ conjecture}

\begin{document}
\begin{abstract}
Assuming the classical integer $abc$ conjecture, we prove the original
conjecture of Rudin. Every arithmetic progression of length $N$ with nonzero
common difference contains $O(\sqrt N)$ squares, with an implied constant
independent of the initial term and common difference. Five representations of one integer as
a sum of two squares give quartic equations via the Pl\"ucker relation.
The coefficients of these equations are bounded polynomially in $N$.
A weighted cut argument then bounds the common difference polynomially
in $N$, at which point a congruence count gives the square-root bound.
\end{abstract}
\maketitle

\section{Introduction}

For integers $a$ and $q\ne0$, let $Q(N;q,a)$ be the number of squares among
$a,a+q,\ldots,a+(N-1)q$, and set $Q(N)=\max_{q\ne0,a}Q(N;q,a)$.
Rudin's conjecture asserts that $Q(N)\ll\sqrt N$ \cite{Rudin,CG}.
This order is sharp, as the progression $1,2,\ldots,N$ shows.

For a positive integer $n$, put $\rad(n)=\prod_{p\mid n}p$.
We use the usual $abc$ inequality in the following form: for a fixed
$\varepsilon>0$ there is a constant $C_\varepsilon$ such that
\begin{equation}\label{eq:abc}
 W\le C_\varepsilon\,\rad(UVW)^{1+\varepsilon}
\end{equation}
for all coprime positive integers $U,V,W$ with $U+V=W$.

\begin{theorem}\label{thm:main}
If \eqref{eq:abc} holds for a single fixed $0<\varepsilon<1/3$, then
$Q(N)\ll_\varepsilon\sqrt N$.
\end{theorem}

In particular, no $N^\varepsilon$ loss occurs in the exponent of the
conclusion. The theorem does not establish the stronger conjecture of
Gonz\'alez-Jim\'enez and Xarles \cite{GJX} that $Q(N)=Q(N;24,1)$ for
$N\ge6$.

\medskip
\noindent\textbf{Related work.}
Bombieri, Granville and Pintz \cite{BGP} obtained the unconditional bound
$Q(N)\ll N^{2/3}(\log N)^{O(1)}$; Bombieri and Zannier \cite{BZ}
improved the exponent to $3/5$. Bourgain and Demeter \cite{BD} obtain
the square-root bound when the common
difference has a bounded number of divisors; Novakovi\'c \cite{Novakovic}
obtains the exponent $1/2+\eta$ for polynomially bounded differences,
for every fixed $\eta>0$. Theorem~\ref{thm:main}
imposes no restriction on the common difference and has no logarithmic
loss, but assumes \eqref{eq:abc}.

\medskip
\noindent\textbf{What the proof adds.}
Cilleruelo and Granville \cite[Section~9]{CG} already proposed using $abc$:
from five square terms they construct a three-term relation whose degree
depends on the gaps between the terms. Their estimate under a strengthening
of $abc$ due to Baker leaves an exponential range for the progression
parameters. More recently, Croot, Mao and Yip
\cite[Theorem~2.5 and Remark~2.6]{CMY} obtained almost sharp conditional
bounds for powers $k\ge4$ and did not treat $k=2,3$ in that theorem.

Here we use five \emph{representations of one sum of two squares}, rather
than five individual square terms. For each quartet the Pl\"ucker relation
gives an equation with three fourth powers. Gaussian factorization and
rounding the prime valuations make the coefficients polynomially bounded
in $N$ and the three fourth-power bases pairwise coprime. Thus $abc$ bounds
each base polynomially: the largest base contributes its fourth power
to the size, whereas the three bases together contribute at most its
third power to the radical. Each base collects prime powers attached to
one balanced partition of the quartet. The five choices of quartet detect
every $2+3$ cut on the five labels; comparison with all ten pairwise
Gaussian gcds then bounds the circle norm, and hence the common difference,
polynomially in $N$. The final congruence count gives $O(\sqrt N)$.

\section{Five representations on a circle}

Reverse the progression if necessary to make $q>0$. We may assume it
contains a square. Write $\gcd(a,q)=b^2g$ with $g$ squarefree and put
$a=b^2g A$, $q=b^2g Q$, where $\gcd(A,Q)=1$.
If $a+qn=x^2$, then $bg$ divides $x$ and $g$ divides $A+Qn$.
Existence of such an index implies $\gcd(g,Q)=1$. The admissible indices
thus lie in one residue class $n=n_0+gk$, and division by $(bg)^2$
gives the progression
\[
 \frac{A+Qn_0}{g}+Qk.
\]
It has length at most $N$, has the same number of squares, and its
initial term is coprime to $Q$. We may therefore assume
$\gcd(a,q)=1$. The term zero contributes at most one square and is
ignored below.

Suppose there are $m>5\sqrt N$ positive squares. The $\binom m2$ unordered
pairs of distinct square terms have at most $2N-1$ index sums.
Some index sum $t$ therefore occurs for at least five pairs.
Write the corresponding representations as
\begin{equation}\label{eq:circle}
 \begin{gathered}
 S=x_i^2+y_i^2=2a+qt,\quad x_i>y_i>0,
 \\
 d_i=\frac{x_i^2-y_i^2}{q}\in\{1,\ldots,N-1\},
 \qquad 1\le i\le5.
 \end{gathered}
\end{equation}
The $d_i$ are distinct and $d_i\equiv t\pmod 2$.
Order $z_i=x_i+\ii y_i$ by increasing argument, so that $d_i>d_j$
when $i<j$.

\section{Gaussian factorizations of the minors}

Put $a_i=qd_i$ and $b_i=2x_i y_i$, so that $z_i^2=a_i+\ii b_i$ and
$a_i^2+b_i^2=S^2$. For $i<j$ write
\[
 D_{ij}=a_i b_j-b_i a_j,\qquad E_{ij}=a_i b_j+b_i a_j.
\]
Both are positive. Let $g_{ij}$ be a Gaussian gcd of $z_i,z_j$, and put
$u=z_i/g_{ij}$, $v=z_j/g_{ij}$. These are coprime and have equal norm,
so $u\bar u=v\bar v$ implies that $v$ is a unit multiple of $\bar u$.
Set
\[
 K'_{ij}=|g_{ij}|^2,\qquad K_{ij}=|u|^2,\qquad K_{ij}K'_{ij}=S .
\]
Taking imaginary parts in $\bar z_i^2 z_j^2$ and $z_i^2 z_j^2$
shows that $(K'_{ij})^2\mid D_{ij}$ and $K_{ij}^2\mid E_{ij}$.

In fact
\begin{equation}\label{eq:J}
 J^-_{ij}=\frac{D_{ij}}{q(K'_{ij})^2},\qquad
 J^+_{ij}=\frac{E_{ij}}{qK_{ij}^2}
\end{equation}
are integers. Indeed, $q$ divides both minors and $\gcd(q,S)\mid2$.
If $q$ is even then $x_i,y_i$ are all odd, so $v_2(S)=1$,
$v_2(K'_{ij})=1$, and $v_2(K_{ij})=0$. Now
$D_{ij}=2q(d_i x_j y_j-d_j x_i y_i)$ and the expression in parentheses
is even since $d_i\equiv d_j\pmod2$. This supplies the additional
factor of four required for $J^-_{ij}$.

The identity $D_{ij}E_{ij}=q^2 S^2(d_i^2-d_j^2)$ gives
\begin{equation}\label{eq:J-bound}
 J^-_{ij}J^+_{ij}=d_i^2-d_j^2,\qquad
 1\le J^\pm_{ij}\le N^2.
\end{equation}
Also $D_{ij}\le a_i b_j\le qNS$. By \eqref{eq:J} we obtain
\begin{equation}\label{eq:K-bound}
 (K'_{ij})^2\le NS,\qquad K_{ij}\ge\sqrt{S/N}.
\end{equation}

\section{A quartic equation from four representations}

Choose any four of the five indices and relabel them in angular order.
The Pl\"ucker relation for the minors of the vectors $(a_i,b_i)$ reads
\[
 D_{12}D_{34}+D_{14}D_{23}=D_{13}D_{24}.
\]
Number the three pairings in the order of these terms. For a pairing
$\nu=(ij,kl)$ define $X_\nu=K'_{ij}K'_{kl}$ and $A_\nu=J^-_{ij}J^-_{kl}$.
After cancelling $q^2$ the relation becomes
\begin{equation}\label{eq:plucker}
 A_1X_1^2+A_2X_2^2=A_3X_3^2,\qquad 1\le A_\nu\le N^4.
\end{equation}

For each odd split prime power $p^e\Vert S$, fix a Gaussian prime $\pi$
above $p$. Let $\alpha_i=v_\pi(z_i)$, so that
$v_{\bar\pi}(z_i)=e-\alpha_i$. Choose $\beta_i\in\{0,1\}$ with
$e\beta_i$ nearest to $\alpha_i$, choosing a tie once for each $p,i$
throughout the proof. Then
\begin{equation}\label{eq:valuation}
 |\alpha_i-e\beta_i|\le\tfrac12 v_p(d_i),\qquad
 v_p(K'_{ij})=e-|\alpha_i-\alpha_j|.
\end{equation}
For the first inequality, $p^{\min(\alpha_i,e-\alpha_i)}$ divides both
$x_i$ and $y_i$, and $p\nmid q$. For the second, take the smaller
valuation at each of $\pi,\bar\pi$ in the gcd.

At the prime $p$, the idealized valuation of $X_\nu$ obtained by
replacing $\alpha_i$ with $e\beta_i$ is
$e(\mathbf{1}_{\beta_i=\beta_j}+\mathbf{1}_{\beta_k=\beta_l})$.
The error is at most $\tfrac12\sum_i v_p(d_i)$ by \eqref{eq:valuation}.
For a $2+2$ split of the labels, the idealized valuations of the
three $X_\nu$ are $2e,0,0$ in some order; the valuation $2e$
belongs to the pairing that matches equal labels. In all other cases
the three idealized valuations agree. Let $P_\nu$ be the product of
the $p^e$ whose labels form a $2+2$ split with equal labels paired by $\nu$.
The three integers $P_\nu$ are pairwise coprime. At inert primes and at
$2$ the actual valuations of all three $X_\nu$ agree: a Gaussian prime
fixed up to units by conjugation has the same valuation in every $z_i$
of norm $S$.

Let $H=\gcd(X_1,X_2,X_3)$ and $D_0=\prod_i d_i\le N^4$. Subtracting
the minimum valuation from each of the three valuations changes the
rounding error by at most another $\tfrac12\sum_i v_p(d_i)$.
Consequently
\[
 \frac{X_\nu}{H}=P_\nu^2 t_\nu
\]
where $t_\nu>0$ is rational and its reduced numerator and denominator
each divide $D_0$. Divide \eqref{eq:plucker} by $H^2$ and multiply by
$D_0^2$. Since $D_0 t_\nu$ is an integer at most $D_0^2$, we obtain
\begin{equation}\label{eq:quartic}
 c_1P_1^4+c_2P_2^4=c_3P_3^4,\qquad
 c_\nu=A_\nu(D_0t_\nu)^2\in\Z,\qquad 1\le c_\nu\le N^{20}.
\end{equation}

Put $M=\max_\nu P_\nu$. The gcd of the three terms in
\eqref{eq:quartic} divides $c_1c_2c_3$, since a prime divides at most
one $P_\nu$. After division by this gcd the largest term is at least
$M^4/N^{60}$, while the radical of the product is at most $N^{60}M^3$.
Applying \eqref{eq:abc} gives
\[
 M^{1-3\varepsilon}\le C_\varepsilon N^{120+60\varepsilon}.
\]
In particular, for each of the five possible quadruples,
\begin{equation}\label{eq:P-bound}
 \log P_\nu\ll_\varepsilon\log(2N).
\end{equation}

\section{Weighted cuts}

Restore all five indices. Each odd split prime power $p^e\Vert S$
defines a cut on the five labels, with weight $e\log p$. For a pair
$i,j$, the total weight of the cuts separating $i$ and $j$ differs
from $\log K_{ij}$ by at most $(\log d_i+\log d_j)/2$, by
\eqref{eq:valuation}. Inert primes and $2$ do not contribute to
$K_{ij}$. Summing over the ten pairs and using \eqref{eq:K-bound}
shows that the sum of the ten separation weights is at least
$5\log S-15\log N$.

Let $T$ be the total weight of the cuts with parts of sizes two and
three. Such a cut separates six pairs; any other cut separates at
most four. Since the total weight is at most $\log S$,
\begin{equation}\label{eq:cuts}
 5\log S-15\log N\le 4\log S+2T.
\end{equation}
Every cut with parts of sizes two and three restricts to a $2+2$ cut
on at least one of the five quadruples. Thus $T$ is at most the sum
of the logarithms in \eqref{eq:P-bound} over all five quadruples
and all three pairings in each. It follows from \eqref{eq:cuts} that
$\log S\ll_\varepsilon\log(2N)$. Since $qd_i<S$, there is a fixed
$B=B(\varepsilon)$ such that $q\ll_\varepsilon N^B$.

\section{Counting roots modulo the common difference}

The positive roots lie in an interval of length at most
$\sqrt{qN}$ and satisfy $x^2\equiv a\pmod q$. Since $\gcd(a,q)=1$,
there are at most two roots modulo each odd prime power dividing $q$
and at most four modulo the power of $2$. The Chinese remainder theorem
therefore gives
\[
 m\ll 2^{\omega(q)}\left(1+\sqrt{N/q}\right),
\]
where $\omega(q)$ is the number of distinct prime divisors of $q$.
If $q\le N$, the bound $2^{\omega(q)}\ll\sqrt q$ gives
$m\ll\sqrt N$. If $N<q\ll_\varepsilon N^B$, fix
$0<\eta<1/(2B)$ and apply $2^{\omega(q)}\ll_\eta q^\eta$.
Both divisor estimates follow by separating the finitely many primes
with $p^\eta<2$ from the rest. This proves $m\ll_\varepsilon\sqrt N$
also in the second range. Together with the case $m\le5\sqrt N$ and
the initial normalization, this proves Theorem~\ref{thm:main}.
\qed

\section*{Lean verification}
The \href{https://github.com/alexlegeartis/RudinConjecture/tree/main/lean}{Lean 4
formalization} follows the same argument at $\varepsilon=1/4$, which is
sufficient to deduce Rudin's conjecture from the classical $abc$ conjecture.
The $abc$ statement is an explicit hypothesis, not an added axiom.
The repository contains a proof crosswalk and a reproducible axiom audit.

\section*{Acknowledgments}
The author posed the problem and asked whether the Riemann hypothesis
could be used to prove Rudin's conjecture. OpenAI's Dot, using GPT-6.2
Astra with Ultra reasoning effort, proposed instead to derive Rudin's
conjecture from the classical \(abc\) conjecture and developed the
conditional proof presented here. At the author's request, Dot also
formalized the same argument in Lean 4.

\bibliographystyle{amsplain}
\bibliography{references}
\end{document}
